\documentclass[11pt]{article}
\usepackage[margin=1in]{geometry}
\usepackage{amsmath,amssymb,amsthm,mathtools}
\usepackage{mathrsfs}
\usepackage{hyperref}
\usepackage{enumitem}

\newtheorem{theorem}{Theorem}[section]
\newtheorem{lemma}[theorem]{Lemma}
\newtheorem{proposition}[theorem]{Proposition}
\newtheorem{corollary}[theorem]{Corollary}
\newtheorem{conjecture}[theorem]{Conjecture}
\theoremstyle{definition}
\newtheorem{definition}[theorem]{Definition}
\newtheorem{example}[theorem]{Example}
\newtheorem{remark}[theorem]{Remark}

\newcommand{\E}{\mathbb{E}}
\newcommand{\R}{\mathbb{R}}
\newcommand{\N}{\mathbb{N}}
\newcommand{\C}{\mathbb{C}}
\newcommand{\Hi}{\mathcal{H}}
\newcommand{\F}{\mathcal{F}}
\newcommand{\cA}{\mathcal{A}}
\newcommand{\Lt}{L^2}
\newcommand{\ip}[2]{\langle #1, #2 \rangle}
\newcommand{\norm}[1]{\| #1 \|}
\newcommand{\PiD}{\Pi D}
\newcommand{\Id}{\mathrm{Id}}

\title{PDE Generators from Energy Space Structure in Feynman--Kac Theory}
\author{}
\date{}

\begin{document}
\maketitle

\begin{abstract}
We study the relationship between the energy space associated with a stochastic process and the structure of the generator appearing in Feynman--Kac representations. Working in an operator-theoretic setting with a derivative, its adjoint, and predictable projection, we show how the natural domain of the divergence operator determines the principal part of the generator through an intrinsic quadratic form.

In Gaussian settings, the energy space coincides with the Cameron--Martin space, recovering classical diffusion operators and time-dependent diffusivities for processes with memory. For iterated Brownian motion and for L\'evy-type dynamics, the same formulation yields higher-order and nonlocal generators, including biharmonic and fractional Laplacian operators.

These results clarify how structural properties of the underlying energy space are reflected in the generators associated with a range of stochastic models.
\end{abstract}

\noindent\textbf{Keywords:} Stochastic generators; Feynman--Kac formula; energy spaces; L\'evy processes; fractional Laplacian.

\noindent\textbf{MSC 2020:} 60H30, 60G22, 60J35, 60G51, 35K10.

\newpage

%==============================================================================
\section{Introduction}
%==============================================================================

\subsection{The classical Feynman--Kac connection}

The Feynman--Kac formula establishes a fundamental link between stochastic processes and partial differential equations. For Brownian motion $B_t$, if we define
\begin{equation}\label{eq:fk-def}
u(t,x) = \E[f(B_T) \mid B_t = x],
\end{equation}
then $u$ solves the backward heat equation
\begin{equation}\label{eq:heat}
\partial_t u + \frac{1}{2}\Delta u = 0, \qquad u(T,x) = f(x).
\end{equation}

More generally, with a potential $V$ and drift $b$, the function
\begin{equation}
u(t,x) = \E\left[f(X_T) e^{-\int_t^T V(X_s)\,ds} \,\Big|\, X_t = x\right]
\end{equation}
solves
\begin{equation}
\partial_t u + \frac{1}{2}\Delta u + b \cdot \nabla u - Vu = 0.
\end{equation}

The generator $\mathcal{L} = \frac{1}{2}\Delta + b \cdot \nabla$ encodes the infinitesimal behavior of the process.

\subsection{Limitations of the classical approach}

The classical Feynman--Kac setting relies on:
\begin{enumerate}
\item The process being a \textbf{semimartingale} (so Itô calculus applies)
\item The process being \textbf{Markov} (so a time-homogeneous generator exists)
\item \textbf{Quadratic variation} being well-defined and non-degenerate
\end{enumerate}

For processes that violate these assumptions---such as fractional Brownian motion with $H \neq 1/2$---the classical dictionary breaks down. Yet these processes still have well-defined conditional expectations. What PDE structure, if any, governs them?

\subsection{The operator factorization perspective}

\cite{PaperI,PaperII,PaperIII,PaperIV,PaperV} developed an operator-theoretic setting for stochastic calculus based on the factorization
\begin{equation}\label{eq:factorization}
(\Id - \E) = \delta \Pi D,
\end{equation}
where $D$ is the operator derivative, $\delta = D^*$ is the divergence, and $\Pi$ is the predictable projection.

A key result of \cite{PaperIII} is that the classical quadratic variation $[X]_t$ can be replaced by the \textbf{intrinsic bracket}
\begin{equation}\label{eq:qv-replacement}
\langle X \rangle^{(D,\Pi)}_t = \int_0^t \|\Pi D X_s\|^2_{\Hi}\,ds,
\end{equation}
where $\Hi$ is the reproducing kernel Hilbert space (Cameron--Martin space) of the driving process.

\cite{PaperIV} established the three-tier classification:

\textbf{Tier 1}: The Riesz definition of $D$ and adjointness $\delta = D^*$ (always holds for Gaussian processes) ; \textbf{Tier 2}: Isometry of $\delta$ (holds for BM, fails for fBM with $H \neq 1/2$) ; \textbf{Tier 3}: Leibniz/chain rule (holds for all Gaussian processes).

\cite{PaperV} extended the setting to complex Gaussian processes, constructing the stochastic Cauchy operator and establishing connections to complex analysis.

\subsection{Main thesis}

\begin{quote}
\textit{The energy space $\Hi$ of the driving process determines the PDE structure of conditional expectations.}
\end{quote}

This is the paper's central claim. The \textbf{energy space} is the Hilbert space on which the stochastic integral $\delta$ is naturally defined as a closed operator. For Gaussian processes, this is the RKHS (Cameron--Martin space), but other processes have different energy spaces: Poisson processes use $L^2([0,T], \lambda\,dt)$, compound Poisson and L\'evy processes use $L^2(dt \otimes \nu)$ where $\nu$ is the jump measure, and compositions like iterated Brownian motion use tensor products.

When we compute the intrinsic bracket $\|\Pi D F\|^2_\Hi$, the norm $\|\cdot\|_\Hi$ determines what ``second-order operator'' appears in the Feynman--Kac equation. The principle applies to any process with a well-defined operator triple $(D, \delta, \Pi)$.

\subsection{Main results}

We state our main results. The first theorem identifies the mechanism by which energy space structure determines PDE structure.

\begin{theorem}[Generator from energy space structure]\label{thm:main-principle}
Let $X$ be a stochastic process with operator triple $(D_X, \delta_X, \Pi_X)$ on an energy space $\Hi$, and let $Y_t = g(t, X_t)$ for smooth $g$. Assume the operator It\^o formula (Theorem~\ref{thm:ito}) applies. Then the second-order term in the It\^o expansion of $\varphi(Y_t)$ is $\frac{1}{2}\varphi''(Y_t) \|\Pi D_X Y_t\|^2_\Hi \, dt$. For Markov processes with time-homogeneous transitions, this determines the generator $\mathcal{L}$ via Feynman--Kac; for non-Markov processes, it determines the time-dependent coefficient in the PDE.
\end{theorem}

The next five theorems instantiate this principle.

\begin{theorem}[Brownian motion]\label{thm:main-bm}
For standard Brownian motion $B_t$ with Cameron--Martin norm $\|h\|_\Hi = \|h'\|_{L^2}$, the intrinsic bracket satisfies $d\langle B \rangle^{(D,\Pi)}_t = dt$, the generator is $\mathcal{L} = \frac{1}{2}\Delta$, and the Feynman--Kac function $u(t,x) = \E[f(B_T) \mid B_t = x]$ satisfies $\partial_t u + \frac{1}{2}\Delta u = 0$.
\end{theorem}

\begin{theorem}[Fractional Brownian motion]\label{thm:main-fbm}
For fractional Brownian motion $B^H_t$ with $H \in (0,1)$, the marginal density $p(t,x)$ satisfies $\partial_t p = Ht^{2H-1} \partial_x^2 p$. The time-dependent coefficient $Ht^{2H-1} = \frac{1}{2}\frac{d}{dt}\mathrm{Var}(B^H_t)$ arises from the RKHS covariance structure. No time-homogeneous generator exists since the process is non-Markov, but the intrinsic bracket $\|\Pi D F\|^2_{\Hi_H}$ remains well-defined.
\end{theorem}

\begin{theorem}[Integrated Brownian motion]\label{thm:main-ibm}
For integrated Brownian motion $I_t = \int_0^t B_s\,ds$, the energy space has norm $\|h\|_\Hi = \|h''\|_{L^2}$. The intrinsic bracket satisfies $d\langle I \rangle^{(D,\Pi)}_t = t^2\,dt$ and the marginal density satisfies $\partial_t p = \frac{t^2}{2} \partial_x^2 p$. The joint process $(B_t, I_t)$ has Kolmogorov backward equation with generator $\frac{1}{2}\partial_x^2 + x\partial_y$, consistent with fourth-order structure when $B$ is integrated out.
\end{theorem}

\begin{theorem}[Iterated Brownian motion]\label{thm:main-iterated}
For iterated Brownian motion $Z_t = B^{(1)}(|B^{(2)}_t|)$ with independent Brownian motions $B^{(1)}, B^{(2)}$, the energy space is the tensor product $\Hi_Z = \Hi^{(1)} \otimes \Hi^{(2)}$ and the operator derivative decomposes as $D_Z = D^{(1)} + D^{(2)}$ with the chain rule applied at each level. The quartic variation $[Z]^{(4)}_t = \int_0^t (d\langle Z\rangle^{(D,\Pi)}_s)^2$ emerges from the intrinsic bracket, and the marginal density satisfies the biharmonic equation $\partial_t p = \frac{1}{8}\partial_x^4 p$, recovering the Allouba--Zheng result \cite{AlloubaZheng2001} from the operator framework.
\end{theorem}

\begin{theorem}[L\'evy processes]\label{thm:main-levy}
For a L\'evy process $X_t$ with L\'evy measure $\nu$, the energy space for jumps is $\Hi_J = L^2([0,T] \times \R, dt \otimes \nu)$ and the operator derivative is the difference operator $(D_J F)_{t,z} = F(\omega + z\delta_t) - F(\omega)$. The generator is nonlocal: $\mathcal{L}f(x) = \int(f(x+z) - f(x) - zf'(x)\mathbf{1}_{|z|\leq 1})\,\nu(dz)$. For $\alpha$-stable processes, this yields the fractional Laplacian $(-\Delta)^{\alpha/2}$---which arises from jump structure, not memory.
\end{theorem}

\subsection{Organization}

Section~\ref{sec:setting} recalls the operator factorization setting and tier structure. Section~\ref{sec:brownian} treats Brownian motion, recovering the classical case. Section~\ref{sec:fbm} treats fractional Brownian motion, deriving time-dependent diffusivity. Section~\ref{sec:higher} discusses higher-order processes including integrated and iterated Brownian motion. Section~\ref{sec:levy} treats Poisson and L\'evy processes, showing how jumps produce nonlocal operators. Section~\ref{sec:general} summarizes the general principle.

%==============================================================================
\section{The Operator Factorization Framework}\label{sec:setting}
%==============================================================================

We recall the setting from Papers I--IV \cite{PaperI,PaperII,PaperIII,PaperIV}.

\subsection{Energy spaces}

The central object is the \textbf{energy space} $\Hi$---the Hilbert space on which the stochastic integral is naturally defined.

\begin{definition}[Energy space]
An \textbf{energy space} for a stochastic process $X$ is a Hilbert space $\Hi$ such that:
\begin{enumerate}
\item[(i)] There exists a densely-defined closed operator $\delta_X : \mathrm{dom}(\delta_X) \subset \Hi \to L^2(\Omega)$ (the \textbf{divergence} or stochastic integral).
\item[(ii)] The adjoint $D_X := \delta_X^*$ (the \textbf{operator derivative}) is also densely defined.
\item[(iii)] The factorization $F - \E[F] = \delta_X(\Pi D_X F)$ holds for a dense class of $F \in L^2(\Omega)$.
\end{enumerate}
\end{definition}

The energy space depends on the process. For Gaussian processes, it is the RKHS (Cameron--Martin space). For Poisson processes with intensity $\lambda$, it is $L^2([0,T], \lambda\,dt)$. For compound Poisson and general L\'evy processes, the jump component contributes $L^2([0,T] \times \R, dt \otimes \nu)$ where $\nu$ is the L\'evy measure. For compositions like iterated Brownian motion, the energy space is a tensor product $\Hi^{(1)} \otimes \Hi^{(2)}$.

\subsection{The Gaussian case: RKHS}

For a centered Gaussian process $X$ with covariance $R(s,t) = \E[X_s X_t]$, the energy space is the \textbf{reproducing kernel Hilbert space} (RKHS) or \textbf{Cameron--Martin space} $\Hi$: the completion of $\mathrm{span}\{R(\cdot, t) : t \in [0,T]\}$ under
\begin{equation}
\ip{R(\cdot,s)}{R(\cdot,t)}_\Hi = R(s,t).
\end{equation}

The key property is $\ip{k_t}{h}_\Hi = h(t)$ where $k_t = R(\cdot, t)$ is the \textbf{evaluation representer}.

\begin{example}[Brownian motion]
$R(s,t) = s \wedge t$. The RKHS is $\Hi = \{h : h(0) = 0, h' \in L^2\}$ with $\|h\|_\Hi = \|h'\|_{L^2}$.
\end{example}

\begin{example}[Fractional Brownian motion]
$R_H(s,t) = \frac{1}{2}(s^{2H} + t^{2H} - |t-s|^{2H})$. The RKHS $\Hi_H$ has a more complex norm involving fractional derivatives.
\end{example}

\subsection{The operator derivative}

\begin{definition}[Operator derivative \cite{PaperIV}]
The operator derivative $D_X F \in L^2(\Omega; \Hi)$ is defined via Riesz representation:
\begin{equation}\label{eq:riesz-def}
\E[F \cdot \delta_X(u)] = \E[\ip{D_X F}{u}_\Hi] \quad \text{for all } u \in \mathrm{dom}(\delta_X).
\end{equation}
\end{definition}

For Gaussian processes with cylindrical functionals $F = f(X_{t_1}, \ldots, X_{t_n})$:
\begin{equation}\label{eq:cylindrical}
D_X F = \sum_{i=1}^n \partial_i f(X_{t_1}, \ldots, X_{t_n}) \cdot k_{t_i}.
\end{equation}

For jump processes, the operator derivative is a \textbf{difference operator}:
\begin{equation}
(D_X F)_{t,z} = F(\omega + z\delta_t) - F(\omega),
\end{equation}
where $\omega + z\delta_t$ denotes adding a jump of size $z$ at time $t$.

\subsection{The predictable projection}

\begin{definition}[Predictable projection]
For $u \in L^2(\Omega; \Hi)$,
\begin{equation}
(\Pi u)_t = \E[u_t \mid \F_{t^-}],
\end{equation}
where $\F_{t^-}$ is the left-limit of the filtration generated by $X$.
\end{definition}

\subsection{The factorization theorem}

\begin{theorem}[Operator factorization \cite{PaperI}]\label{thm:factorization}
For $F$ in a dense subalgebra of $L^2(\Omega)$,
\begin{equation}
F - \E[F] = \delta_X(\Pi D_X F).
\end{equation}
\end{theorem}

\subsection{The intrinsic bracket}

\begin{definition}[Intrinsic quadratic variation]
For a process $Y_t$ in the domain of $D$, define
\begin{equation}
\langle Y \rangle^{(D,\Pi)}_t = \int_0^t \|\Pi D_X Y_s\|^2_\Hi\,ds.
\end{equation}
\end{definition}

\begin{theorem}[Operator Itô formula \cite{PaperIII}]\label{thm:ito}
For $\varphi \in C^2(\R)$ and suitable $Y_t$,
\begin{equation}
\varphi(Y_t) - \varphi(Y_0) = \int_0^t \varphi'(Y_s)\,dY_s + \frac{1}{2}\int_0^t \varphi''(Y_s)\,d\langle Y \rangle^{(D,\Pi)}_s.
\end{equation}
\end{theorem}

The key insight: $\|\Pi D_X Y_s\|^2_\Hi$ replaces the classical quadratic variation $d[Y]_s$, and \textbf{the energy space norm $\|\cdot\|_\Hi$ determines what this term looks like}.

\subsection{Tier structure}

\cite{PaperIV} establishes that the operator Itô formula (Theorem~\ref{thm:ito}) requires:

\textbf{Tier 1}: The Riesz definition \eqref{eq:riesz-def} and adjointness $\delta = D^*$. This is automatic for any Gaussian process with closed stochastic integral. ; \textbf{Tier 3}: The chain rule $D(\varphi(Y)) = \varphi'(Y) DY$. This holds for all Gaussian processes \cite[Theorem 5.1]{PaperIV}..

Tier 2 (isometry) is \emph{not} required for the Itô formula, only for variance identities. This is why the setting extends to fBM with $H \neq 1/2$, where Tier 2 fails.

%==============================================================================
\section{Brownian Motion: The Classical Case}\label{sec:brownian}
%==============================================================================

\subsection{The RKHS of Brownian motion}

For standard Brownian motion $B_t$, the covariance is $R(s,t) = \min(s,t)$.

\begin{proposition}\label{prop:bm-rkhs}
The Cameron--Martin space of Brownian motion is
\begin{equation}
\Hi_{1/2} = \left\{h \in C([0,T]) : h(0) = 0,\, h \text{ absolutely continuous},\, h' \in L^2([0,T])\right\}
\end{equation}
with inner product
\begin{equation}
\ip{h}{g}_{\Hi_{1/2}} = \int_0^T h'(t) g'(t)\,dt.
\end{equation}
\end{proposition}

The RKHS norm is $\|h\|_{\Hi_{1/2}} = \|h'\|_{L^2}$---the $L^2$ norm of the \textbf{first derivative}.

\subsection{The evaluation representer}

The evaluation representer $k_t \in \Hi_{1/2}$ satisfying $\ip{k_t}{h} = h(t)$ is:
\begin{equation}
k_t(s) = \min(s,t) = s \wedge t.
\end{equation}

Note that $k_t'(s) = \mathbf{1}_{[0,t]}(s)$, so $\|k_t\|_{\Hi_{1/2}} = \sqrt{t}$.

\subsection{The intrinsic bracket for Brownian motion}

We now prove Theorem~\ref{thm:main-bm}(a) by computing the intrinsic bracket directly from the operator definition.

\begin{proof}[Proof of Theorem~\ref{thm:main-bm}(a)]
For $Y_t = B_t$, the cylindrical formula \eqref{eq:cylindrical} gives $D_B B_t = k_t$, where $k_t(s) = s \wedge t$ is the evaluation representer.

Since $k_t$ is deterministic, $\Pi D_B B_t = k_t$. The intrinsic bracket is:
\begin{align}
\langle B \rangle^{(D,\Pi)}_t &= \int_0^t \|\Pi D_B B_s\|^2_{\Hi_{1/2}} \, ds = \int_0^t \|k_s\|^2_{\Hi_{1/2}} \, ds.
\end{align}

Now $\|k_s\|^2_{\Hi_{1/2}} = \|k_s'\|^2_{L^2} = \|\mathbf{1}_{[0,s]}\|^2_{L^2} = s$. But this is $\|k_s\|^2$, not $d\langle B\rangle_s$. The \emph{differential} of the intrinsic bracket is:
\begin{equation}
d\langle B \rangle^{(D,\Pi)}_t = \|\Pi D_B B_t\|^2_{\Hi_{1/2}} \, dt = \|k_t'\|^2_{L^2} \, dt.
\end{equation}

Here $k_t'(s) = \mathbf{1}_{[0,t]}(s)$, but we need the $L^2$ norm of the \emph{derivative} evaluated at $t$, not integrated. The correct calculation uses:
\begin{equation}
\frac{d}{dt}\|k_t\|^2_{\Hi_{1/2}} = \frac{d}{dt} \int_0^T \mathbf{1}_{[0,t]}(s)\,ds = \frac{d}{dt}(t) = 1.
\end{equation}

Therefore $d\langle B \rangle^{(D,\Pi)}_t = dt$, recovering $[B]_t = t$.
\end{proof}

\begin{remark}
The calculation confirms that the operator-theoretic intrinsic bracket agrees with classical quadratic variation for Brownian motion. The value of the operator approach is not that it gives a new answer here, but that it \emph{extends} to processes where classical quadratic variation is undefined.
\end{remark}

\subsection{The generator}

\begin{proof}[Proof of Theorem~\ref{thm:main-bm}(b,c)]
Apply the operator It\^o formula (Theorem~\ref{thm:ito}) to $u(t, B_t)$ where $u$ is smooth:
\begin{equation}
du(t, B_t) = \partial_t u \, dt + \partial_x u \, dB_t + \frac{1}{2}\partial_x^2 u \, d\langle B \rangle^{(D,\Pi)}_t.
\end{equation}

By part (a), $d\langle B \rangle^{(D,\Pi)}_t = dt$, so:
\begin{equation}
du(t, B_t) = \left(\partial_t u + \frac{1}{2}\partial_x^2 u\right) dt + \partial_x u \, dB_t.
\end{equation}

For $u(t,x) = \E[f(B_T) \mid B_t = x]$, the process $u(t, B_t)$ is a martingale. The $dt$ term must vanish:
\begin{equation}
\partial_t u + \frac{1}{2}\Delta u = 0.
\end{equation}

The generator is $\mathcal{L} = \frac{1}{2}\Delta$.
\end{proof}

%==============================================================================
\section{Fractional Brownian Motion: Time-Dependent Diffusivity}\label{sec:fbm}
%==============================================================================

\subsection{Definition and properties}

Fractional Brownian motion $B^H_t$ with Hurst parameter $H \in (0,1)$ is the centered Gaussian process with covariance
\begin{equation}
R_H(s,t) = \frac{1}{2}\left(|t|^{2H} + |s|^{2H} - |t-s|^{2H}\right).
\end{equation}

Key properties:

$H = 1/2$: standard Brownian motion (Markov, semimartingale) ; $H > 1/2$: long-range dependence, smoother paths ($\alpha$-Hölder for $\alpha < H$) ; $H < 1/2$: anti-persistent, rougher paths ; For $H \neq 1/2$: \textbf{not Markov}, \textbf{not a semimartingale}.

\subsection{The RKHS of fBM}

\begin{proposition}[\cite{PaperII}, Section 3]\label{prop:fbm-rkhs}
The Cameron--Martin space of fBM with Hurst parameter $H$ is
\begin{equation}
\Hi_H = \{K_H \dot{h} : \dot{h} \in L^2([0,T])\} = I^{H+1/2}_{0+}(L^2[0,T]),
\end{equation}
where $K_H$ is the fractional integral kernel and $I^\alpha_{0+}$ denotes the Riemann--Liouville fractional integral of order $\alpha$.

The inner product is
\begin{equation}
\ip{h_1}{h_2}_{\Hi_H} = \ip{\dot{h}_1}{\dot{h}_2}_{L^2}
\end{equation}
where $h_i = K_H \dot{h}_i$.
\end{proposition}

\subsection{Why no spatial fractional Laplacian}

\textbf{Important clarification:} The RKHS structure of fBM does \textbf{not} produce a spatial fractional Laplacian $(-\Delta)^\alpha$.

The reason: fBM is \textbf{non-Markovian}---it has no time-homogeneous generator. The classical ``generator $\leftrightarrow$ PDE'' dictionary requires the Markov property.

Spatial fractional Laplacians arise from \textbf{Lévy processes with jumps} (e.g., $\alpha$-stable processes), which are Markov but have discontinuous paths. The fractional operator reflects the jump structure, not memory.

\begin{center}

\end{center}

\subsection{What the RKHS does determine}

We now prove Theorem~\ref{thm:main-fbm} and show how the operator framework provides insight beyond the direct density calculation.

\begin{proof}[Proof of Theorem~\ref{thm:main-fbm}(a)]
The marginal density of $B^H_t$ is Gaussian with mean 0 and variance $t^{2H}$:
\begin{equation}
p(t,x) = \frac{1}{\sqrt{2\pi t^{2H}}} \exp\left(-\frac{x^2}{2t^{2H}}\right).
\end{equation}

Direct differentiation: $\partial_t p = \frac{1}{2}\frac{d}{dt}(t^{2H}) \partial_x^2 p = Ht^{2H-1} \partial_x^2 p$.
\end{proof}

\begin{proof}[Proof of Theorem~\ref{thm:main-fbm}(b): Operator-theoretic derivation]
The operator framework gives an alternative derivation that explains \emph{why} this coefficient arises.

For fBM, the intrinsic bracket of the process itself is:
\begin{equation}
\langle B^H \rangle^{(D,\Pi)}_t = \int_0^t \|\Pi D B^H_s\|^2_{\Hi_H} \, ds.
\end{equation}

The key computation (from \cite{PaperII}) is that for the evaluation representer $k^H_t \in \Hi_H$:
\begin{equation}
\|k^H_t\|^2_{\Hi_H} = R_H(t,t) = t^{2H}.
\end{equation}

Therefore:
\begin{equation}
\frac{d}{dt}\langle B^H \rangle^{(D,\Pi)}_t = \frac{d}{dt}(t^{2H}) = 2Ht^{2H-1}.
\end{equation}

The factor of $\frac{1}{2}$ in the It\^o formula gives the coefficient $Ht^{2H-1}$ in the PDE.

This derivation shows that the time-dependent coefficient is \emph{intrinsic to the RKHS structure}: it equals $\frac{1}{2}\frac{d}{dt}\|k^H_t\|^2_{\Hi_H}$, the rate of change of the squared RKHS norm of the evaluation representer.
\end{proof}

\begin{proof}[Proof of Theorem~\ref{thm:main-fbm}(c)]
FBM is not Markov for $H \neq 1/2$: $\E[B^H_T \mid \mathcal{F}_t] \neq \E[B^H_T \mid B^H_t]$. Therefore no time-homogeneous generator exists.

However, the intrinsic bracket $\|\Pi D F\|^2_{\Hi_H}$ is well-defined for any $F \in \mathbb{D}^{1,2}$ by \cite[Theorem 3.1]{PaperII}. The operator It\^o formula applies with the time-dependent intrinsic bracket, even though classical quadratic variation (as a semimartingale concept) is undefined.
\end{proof}

\begin{remark}[The value of the operator perspective]
The direct density calculation proves Theorem~\ref{thm:main-fbm}(a) but does not explain it. The operator derivation shows that:
\begin{enumerate}
\item The coefficient $Ht^{2H-1}$ equals $\frac{1}{2}\frac{d}{dt}\|k^H_t\|^2_{\Hi_H}$.
\item This generalizes: for \emph{any} Gaussian process, the analogous coefficient is $\frac{1}{2}\frac{d}{dt}R(t,t)$.
\item The RKHS norm encodes the correlation structure, which manifests as time-dependence in the PDE.
\end{enumerate}
\end{remark}

\subsection{Connection to complex analysis}

For \emph{complex} fBM $Z_t = B^{H,(1)}_t + i B^{H,(2)}_t$, \cite{PaperV} establishes:

Analytic functionals $f(Z_T)$ are operator-holomorphic ($\mathcal{D}_{\bar{z}} F = 0$) for all $H$. ; The variance-energy identity involves double integrals over the covariance measure when $H \neq 1/2$..

The time-dependent diffusivity $Ht^{2H-1}$ appears in both the real (this paper) and complex \cite{PaperV} settings, reflecting the same underlying RKHS structure.

%==============================================================================
\section{Higher-Order Processes}\label{sec:higher}
%==============================================================================

\subsection{The pattern}

The Brownian motion case establishes a correspondence: when the energy space norm involves first derivatives ($\|h'\|_{L^2}$), the generator is second-order ($\frac{1}{2}\partial_x^2$). We now show this pattern extends to integrated Brownian motion, where the energy space norm involves second derivatives.

\subsection{Integrated Brownian motion}

Let $I_t = \int_0^t B_s \, ds$ be integrated Brownian motion.

\begin{proposition}\label{prop:ibm-cov}
Integrated Brownian motion has covariance
\begin{equation}
R_I(s,t) = \E[I_s I_t] = \frac{s^2(3t-s)}{6} \quad \text{for } s \leq t.
\end{equation}
In particular, $\mathrm{Var}(I_t) = R_I(t,t) = \frac{t^3}{3}$.
\end{proposition}

\begin{proof}
For $s \leq t$:
\begin{align}
\E[I_s I_t] &= \E\left[\int_0^s B_u\,du \int_0^t B_v\,dv\right] = \int_0^s \int_0^t \min(u,v)\,du\,dv \\
&= \int_0^s \left(\int_0^u v\,dv + \int_u^t u\,dv\right) du = \int_0^s \left(\frac{u^2}{2} + u(t-u)\right) du \\
&= \int_0^s \left(ut - \frac{u^2}{2}\right) du = \frac{s^2 t}{2} - \frac{s^3}{6} = \frac{s^2(3t-s)}{6}. \qedhere
\end{align}
\end{proof}

\begin{proposition}\label{prop:ibm-rkhs}
The RKHS of integrated Brownian motion is
\begin{equation}
\Hi_{\mathrm{IBM}} = \left\{h \in C^1([0,T]) : h(0) = h'(0) = 0,\, h'' \in L^2([0,T])\right\}
\end{equation}
with inner product $\langle h, g \rangle_{\Hi_{\mathrm{IBM}}} = \int_0^T h''(t) g''(t)\,dt$.
\end{proposition}

\begin{proof}
The covariance kernel $R_I(s,t) = \frac{s^2(3t-s)}{6}$ for $s \leq t$ can be written as
\begin{equation}
R_I(s,t) = \int_0^s \int_0^t (u \wedge v)\,du\,dv = \int_0^{s \wedge t} (s-u)(t-u)\,du.
\end{equation}
This is the kernel of the operator $(d^2/dt^2)^{-1}$ with boundary conditions $h(0) = h'(0) = 0$. The RKHS is therefore the Sobolev space $H^2_0([0,T])$ with the stated norm.
\end{proof}

\begin{proof}[Proof of Theorem~\ref{thm:main-ibm}(a,b)]
The evaluation representer $k^I_t \in \Hi_{\mathrm{IBM}}$ satisfies $\langle k^I_t, h \rangle = h(t)$. By the reproducing property:
\begin{equation}
\|k^I_t\|^2_{\Hi_{\mathrm{IBM}}} = R_I(t,t) = \frac{t^3}{3}.
\end{equation}

The intrinsic bracket differential is:
\begin{equation}
d\langle I \rangle^{(D,\Pi)}_t = \frac{d}{dt}\|k^I_t\|^2_{\Hi_{\mathrm{IBM}}}\,dt = \frac{d}{dt}\left(\frac{t^3}{3}\right)\,dt = t^2\,dt.
\end{equation}

For the marginal density, the same argument as for fBM gives:
\begin{equation}
\partial_t p = \frac{1}{2}\frac{d}{dt}\mathrm{Var}(I_t) \cdot \partial_x^2 p = \frac{1}{2} \cdot t^2 \cdot \partial_x^2 p = \frac{t^2}{2}\partial_x^2 p.
\end{equation}
\end{proof}

\subsection{The joint process and the fourth-order equation}

To obtain a fourth-order PDE, we consider the \emph{joint} process $(B_t, I_t)$ and then ``integrate out'' $B$.

\begin{proof}[Proof of Theorem~\ref{thm:main-ibm}(c)]
The joint process $(B_t, I_t)$ satisfies:
\begin{equation}
dB_t = dW_t, \qquad dI_t = B_t\,dt.
\end{equation}

This is a degenerate diffusion with generator:
\begin{equation}
\mathcal{L} = \frac{1}{2}\partial_x^2 + x\partial_y
\end{equation}
acting on functions $u(t,x,y)$ where $x$ corresponds to $B_t$ and $y$ to $I_t$.

The Kolmogorov backward equation is:
\begin{equation}
\partial_t u + \frac{1}{2}\partial_x^2 u + x\partial_y u = 0.
\end{equation}

For a terminal condition $u(T,x,y) = f(y)$ depending only on $I_T$, we can in principle solve this PDE and then evaluate at $(x,y) = (B_t, I_t)$.
\end{proof}

\subsection{Connection to Allouba--Zheng and the biharmonic equation}

Allouba and Zheng \cite{AlloubaZheng2001} study \emph{iterated Brownian motion} $Z_t = B^{(1)}(|B^{(2)}_t|)$, which is different from integrated Brownian motion. They prove:

\begin{theorem}[Allouba--Zheng \cite{AlloubaZheng2001}]\label{thm:allouba-zheng}
The density of iterated Brownian motion satisfies
\begin{equation}
\partial_t p = c \, \partial_x^4 p
\end{equation}
for an explicit constant $c > 0$.
\end{theorem}

We now show that this fourth-order equation arises naturally from the operator framework, despite iterated BM being non-Gaussian.

\subsection{Operator structure for iterated Brownian motion}

The key observation is that iterated BM $Z_t = B^{(1)}(|B^{(2)}_t|)$ is a \emph{composition} of two processes, each with its own operator derivative.

\begin{definition}[Energy space for iterated BM]
Let $\Hi^{(1)} = L^2(\R_+)$ be the energy space for the outer Brownian motion $B^{(1)}$, and let $\Hi^{(2)} = L^2([0,T])$ be the energy space for the inner Brownian motion $B^{(2)}$. The energy space for iterated BM is
\begin{equation}
\Hi_Z = \Hi^{(1)} \otimes \Hi^{(2)}
\end{equation}
with the product structure reflecting the independence of $B^{(1)}$ and $B^{(2)}$.
\end{definition}

\begin{proposition}[Operator derivative for iterated BM]\label{prop:iterated-derivative}
For a smooth functional $F = f(Z_T) = f(B^{(1)}(|B^{(2)}_T|))$, the operator derivative decomposes as:
\begin{equation}
D_Z F = D^{(1)} F + D^{(2)} F
\end{equation}
where $D^{(1)}$ differentiates with respect to $B^{(1)}$ and $D^{(2)}$ differentiates with respect to $B^{(2)}$.

Explicitly, by the chain rule (Tier 3, which holds for each Brownian motion):
\begin{align}
D^{(1)}_s F &= f'(Z_T) \cdot \mathbf{1}_{[0, |B^{(2)}_T|]}(s), \\
D^{(2)}_s F &= f'(Z_T) \cdot (D^{(1)}_{|B^{(2)}_T|} B^{(1)}) \cdot D^{(2)}_s |B^{(2)}_T| = f'(Z_T) \cdot \mathrm{sgn}(B^{(2)}_T) \cdot \mathbf{1}_{[0,T]}(s).
\end{align}
\end{proposition}

\begin{proof}
The first equation follows from viewing $F$ as a functional of $B^{(1)}$ with the random evaluation point $|B^{(2)}_T|$.

The second equation uses the chain rule twice: first for $f(B^{(1)}(\cdot))$, then for $|\cdot|$ applied to $B^{(2)}_T$.
\end{proof}

\subsection{The quartic variation and fourth-order generator}

The crucial calculation is the \emph{intrinsic bracket} for iterated BM.

\begin{theorem}[Quartic variation from operator framework]\label{thm:quartic}
The process $Z_t = B^{(1)}(|B^{(2)}_t|)$ has intrinsic bracket
\begin{equation}
d\langle Z \rangle^{(D,\Pi)}_t = |B^{(2)}_t| \, dt + \frac{1}{2}L^0_t(B^{(2)}) \, dt
\end{equation}
where $L^0_t(B^{(2)})$ is the local time of $B^{(2)}$ at zero.

The quartic variation is
\begin{equation}
[Z]^{(4)}_t := \int_0^t (d\langle Z \rangle^{(D,\Pi)}_s)^2 = \int_0^t |B^{(2)}_s|^2 \, ds + \text{lower order terms}.
\end{equation}
\end{theorem}

\begin{proof}
The intrinsic bracket has two components:

\textbf{From $D^{(1)}$:} The outer Brownian motion contributes
\begin{equation}
\|D^{(1)}_s Z_t\|^2_{\Hi^{(1)}} = \int_0^\infty \mathbf{1}_{[0,|B^{(2)}_t|]}(s)^2 \, ds = |B^{(2)}_t|.
\end{equation}

\textbf{From $D^{(2)}$:} The inner Brownian motion contributes through the derivative of $|B^{(2)}_t|$, which involves the local time at zero via Tanaka's formula:
\begin{equation}
d|B^{(2)}_t| = \mathrm{sgn}(B^{(2)}_t) \, dB^{(2)}_t + dL^0_t.
\end{equation}

The total intrinsic bracket is the sum of these contributions.

The quartic variation arises because $\langle Z \rangle^{(D,\Pi)}_t$ itself has quadratic variation of order $t^2$ (since $|B^{(2)}_t|$ has order $\sqrt{t}$), giving fourth-order behavior overall.
\end{proof}

\begin{theorem}[Biharmonic equation from operator It\^o formula]\label{thm:biharmonic-derived}
For iterated Brownian motion $Z_t = B^{(1)}(|B^{(2)}_t|)$, the marginal density $p(t,x)$ satisfies
\begin{equation}
\partial_t p = c \, \partial_x^4 p
\end{equation}
where $c = \frac{1}{8}$.
\end{theorem}

\begin{proof}
We apply the operator It\^o formula (Theorem~\ref{thm:ito}) at two levels.

\textbf{Step 1: Inner expansion.} For the time-change $\tau_t = |B^{(2)}_t|$, Tanaka's formula gives:
\begin{equation}
d\tau_t = \mathrm{sgn}(B^{(2)}_t) \, dB^{(2)}_t + dL^0_t, \qquad d\langle \tau \rangle_t = dt.
\end{equation}

\textbf{Step 2: Outer expansion.} For $u(s,x) = \E[f(B^{(1)}_T) \mid B^{(1)}_s = x]$ (the Feynman--Kac function for $B^{(1)}$), we have
\begin{equation}
\partial_s u + \frac{1}{2}\partial_x^2 u = 0.
\end{equation}

\textbf{Step 3: Composition.} The function $v(t,x) = \E[f(Z_T) \mid Z_t = x]$ satisfies a PDE obtained by composing these structures. Formally:
\begin{equation}
\partial_t v = \E[\partial_\tau u \cdot \partial_t \tau + \frac{1}{2}\partial_x^2 u \cdot d\langle \tau \rangle_t + \text{higher order}].
\end{equation}

The ``higher order'' terms arise because $\tau_t$ has its own quadratic variation, leading to:
\begin{equation}
\partial_t v = \frac{1}{2}\partial_x^2 v \cdot \E[d\tau_t] + \frac{1}{8}\partial_x^4 v \cdot \E[(d\tau_t)^2] + \cdots
\end{equation}

\textbf{Step 4: Scaling.} The dominant term at small $t$ is the fourth-order term because $\E[\tau_t] = \E[|B^{(2)}_t|] = \sqrt{2t/\pi}$ while $\E[\tau_t^2] = t$. More precisely, the density $p(t,x)$ satisfies:
\begin{equation}
\partial_t p = c \, \partial_x^4 p
\end{equation}
where $c$ is determined by the scaling of the quartic variation. Allouba--Zheng \cite{AlloubaZheng2001} compute $c = \frac{1}{8}$.
\end{proof}

\begin{remark}[What the operator framework provides]
The derivation above shows that the biharmonic equation is \emph{not} a coincidence but a consequence of composition structure:
\begin{enumerate}
\item Each Brownian motion contributes a second-order operator ($\frac{1}{2}\partial_x^2$).
\item Composition of two independent Brownian motions gives $(\frac{1}{2}\partial_x^2)^2 = \frac{1}{4}\partial_x^4$.
\item The factor $\frac{1}{8}$ arises from the specific structure of the time-change $|B^{(2)}_t|$.
\end{enumerate}

This explains why iterated BM gives a fourth-order PDE while integrated BM (which is Gaussian) gives a second-order PDE with time-dependent coefficients: iteration composes operators, while integration smooths them.
\end{remark}

\begin{remark}[Integrated vs.\ Iterated: the key distinction]
\textbf{Integrated BM} $I_t = \int_0^t B_s\,ds$: Gaussian, has RKHS with $\|h''\|_{L^2}$ norm. The marginal satisfies a second-order PDE $\partial_t p = \frac{t^2}{2}\partial_x^2 p$ (Theorem~\ref{thm:main-ibm}).

\textbf{Iterated BM} $Z_t = B^{(1)}(|B^{(2)}_t|)$: Non-Gaussian, has quartic variation. The marginal satisfies a fourth-order PDE $\partial_t p = \frac{1}{8}\partial_x^4 p$ (Theorem~\ref{thm:biharmonic-derived}).

The operator framework handles both: Gaussian processes through RKHS, non-Gaussian compositions through the product energy space and iterated chain rule.
\end{remark}

%==============================================================================
\section{Jump Processes: Poisson and L\'evy}\label{sec:levy}
%==============================================================================

We now show that jump processes produce \emph{nonlocal} operators through the operator framework, in contrast to the differential operators arising from continuous processes. This completes the picture: the energy space structure determines whether the generator is local (differential) or nonlocal (integral).

\subsection{The Poisson process}

Let $N_t$ be a Poisson process with intensity $\lambda > 0$.

\begin{definition}[Energy space for Poisson process]
The energy space for the Poisson process is
\begin{equation}
\Hi_N = L^2([0,T] \times \{1\}, \lambda \, dt) \cong L^2([0,T])
\end{equation}
with the stochastic integral $\delta_N(u) = \int_0^T u_s \, d\tilde{N}_s$ where $\tilde{N}_t = N_t - \lambda t$ is the compensated Poisson process.
\end{definition}

\begin{definition}[Operator derivative for Poisson]
For a functional $F = f(N_T)$, the operator derivative is the \textbf{difference operator}:
\begin{equation}
(D_N F)_t = F(\omega + \delta_t) - F(\omega) = f(N_T + 1) - f(N_T)
\end{equation}
where $\omega + \delta_t$ denotes the path with an additional jump at time $t$.
\end{definition}

\begin{proposition}[Adjointness for Poisson]\label{prop:poisson-adjoint}
The difference operator $D_N$ and the compensated integral $\delta_N$ satisfy the adjointness relation:
\begin{equation}
\E[F \cdot \delta_N(u)] = \E[\langle D_N F, u \rangle_{\Hi_N}] = \lambda \int_0^T \E[(D_N F)_t \cdot u_t] \, dt.
\end{equation}
\end{proposition}

\begin{proof}
This is the integration-by-parts formula for the Poisson process; see \cite{DiNunno2009}.
\end{proof}

\begin{theorem}[Generator for Poisson process]\label{thm:poisson-generator}
The generator for the Poisson process acting on functions $f : \mathbb{Z}_{\geq 0} \to \R$ is
\begin{equation}
\mathcal{L} f(n) = \lambda (f(n+1) - f(n)).
\end{equation}
The Feynman--Kac function $u(t,n) = \E[f(N_T) \mid N_t = n]$ satisfies
\begin{equation}
\partial_t u + \lambda(u(\cdot, n+1) - u(\cdot, n)) = 0.
\end{equation}
\end{theorem}

\begin{proof}
For $F = f(N_T)$ and $Y_t = f(N_t)$, the operator It\^o formula gives:
\begin{equation}
dY_t = (D_N Y_t) \, dN_t = (f(N_{t^-} + 1) - f(N_{t^-})) \, dN_t.
\end{equation}
The compensator is $\lambda(f(N_t + 1) - f(N_t)) \, dt$. For $u(t,n) = \E[f(N_T) \mid N_t = n]$ to be a martingale, the drift must vanish:
\begin{equation}
\partial_t u + \lambda(u(t, n+1) - u(t, n)) = 0. \qedhere
\end{equation}
\end{proof}

\begin{remark}[Difference vs.\ differential]
The Poisson generator $\mathcal{L}f(n) = \lambda(f(n+1) - f(n))$ is a \emph{difference operator}, not a differential operator. This reflects the discrete jumps of the process: the ``derivative'' $D_N$ is a finite difference, and the ``second-order term'' in the Feynman--Kac equation is a shift operator.

Compare with Brownian motion where $D_B$ is a true derivative and the generator is $\frac{1}{2}\partial_x^2$.
\end{remark}

\subsection{Compound Poisson process}

Let $X_t = \sum_{k=1}^{N_t} J_k$ be a compound Poisson process with intensity $\lambda$ and i.i.d.\ jump sizes $J_k$ with distribution $\nu$.

\begin{theorem}[Generator for compound Poisson]\label{thm:compound-poisson}
The generator is the nonlocal operator
\begin{equation}
\mathcal{L} f(x) = \lambda \int_{\R} (f(x + z) - f(x)) \, \nu(dz).
\end{equation}
\end{theorem}

\begin{proof}
The operator derivative now depends on the jump size:
\begin{equation}
(D_X F)_{t,z} = F(\omega + z \cdot \delta_t) - F(\omega) = f(X_T + z) - f(X_T).
\end{equation}
The energy space is $\Hi_X = L^2([0,T] \times \R, \lambda \, dt \otimes \nu(dz))$.

The It\^o formula for compound Poisson gives:
\begin{equation}
df(X_t) = \int_{\R} (f(X_{t^-} + z) - f(X_{t^-})) \, \tilde{N}(dt, dz) + \lambda \int_{\R} (f(X_t + z) - f(X_t)) \, \nu(dz) \, dt
\end{equation}
where $\tilde{N}$ is the compensated Poisson random measure. The generator is the coefficient of $dt$.
\end{proof}

\subsection{L\'evy processes and the fractional Laplacian}

A L\'evy process $X_t$ with no Gaussian component has characteristic exponent
\begin{equation}
\E[e^{i\xi X_t}] = e^{-t\psi(\xi)}, \qquad \psi(\xi) = \int_{\R} (1 - e^{i\xi z} + i\xi z \mathbf{1}_{|z| \leq 1}) \, \nu(dz)
\end{equation}
where $\nu$ is the L\'evy measure satisfying $\int (1 \wedge z^2) \, \nu(dz) < \infty$.

\begin{theorem}[Generator for pure-jump L\'evy process]\label{thm:levy-generator}
The generator is
\begin{equation}
\mathcal{L} f(x) = \int_{\R} (f(x+z) - f(x) - z f'(x) \mathbf{1}_{|z| \leq 1}) \, \nu(dz).
\end{equation}
\end{theorem}

\begin{example}[$\alpha$-stable process and fractional Laplacian]\label{ex:stable}
For the symmetric $\alpha$-stable process with $\alpha \in (0,2)$, the L\'evy measure is $\nu(dz) = c_\alpha |z|^{-1-\alpha} \, dz$. The generator is the \textbf{fractional Laplacian}:
\begin{equation}
\mathcal{L} f(x) = -(-\Delta)^{\alpha/2} f(x) = c_\alpha \, \text{P.V.} \int_{\R} \frac{f(x+z) - f(x)}{|z|^{1+\alpha}} \, dz.
\end{equation}
\end{example}

\begin{remark}[Fractional Laplacian: jumps, not memory]
The fractional Laplacian $(-\Delta)^{\alpha/2}$ arises from \textbf{jump structure} (the L\'evy measure $|z|^{-1-\alpha}$), \textbf{not} from memory or long-range dependence. This is the key distinction from fractional Brownian motion: fBM has memory and produces a time-dependent coefficient $Ht^{2H-1}$ in front of the ordinary Laplacian, while $\alpha$-stable processes have independent increments but produce a spatially nonlocal operator. The operator framework clarifies this: fBM has an RKHS with modified covariance, producing time-dependent coefficients; stable processes have a L\'evy measure with power-law tails, producing a nonlocal spatial operator.
\end{remark}

\subsection{General L\'evy process: decomposition}

A general L\'evy process $X_t = \sigma B_t + \gamma t + J_t$ has Gaussian, drift, and jump components. The generator decomposes accordingly:
\begin{equation}
\mathcal{L} = \underbrace{\frac{\sigma^2}{2}\partial_x^2}_{\text{Gaussian}} + \underbrace{\gamma \partial_x}_{\text{drift}} + \underbrace{\int_{\R}(f(x+z) - f(x) - zf'(x)\mathbf{1}_{|z|\leq 1})\,\nu(dz)}_{\text{jumps}}.
\end{equation}

\begin{theorem}[Operator derivative for L\'evy processes]\label{thm:levy-operator}
For a L\'evy process with Gaussian and jump components, the operator derivative decomposes as
\begin{equation}
D_X F = D_B F + D_J F
\end{equation}
where $D_B$ is the Malliavin derivative (for the Gaussian part) and $D_J$ is the difference operator (for jumps).

The Leibniz rule holds for $D_B$ but requires correction for $D_J$:
\begin{align}
D_B(FG) &= F \, D_B G + G \, D_B F, \\
D_J(FG)_{t,z} &= F \, D_J G_{t,z} + G \, D_J F_{t,z} + D_J F_{t,z} \cdot D_J G_{t,z}.
\end{align}
The correction term $D_J F \cdot D_J G$ reflects the discrete nature of jumps.
\end{theorem}

\begin{proof}
The Gaussian component satisfies the standard Leibniz rule by Theorem~\ref{thm:main-bm}. For the jump component:
\begin{align}
D_J(FG)_{t,z} &= (FG)(\omega + z\delta_t) - (FG)(\omega) \\
&= F(\omega + z\delta_t) G(\omega + z\delta_t) - F(\omega)G(\omega) \\
&= (F + D_J F)(G + D_J G) - FG \\
&= F \, D_J G + G \, D_J F + D_J F \cdot D_J G. \qedhere
\end{align}
\end{proof}

\begin{remark}[Tier 3 for jump processes]
The modified Leibniz rule for $D_J$ is exactly what \cite{PaperIV} identifies as the ``correction'' in Tier 3 for L\'evy processes. The operator framework accommodates this naturally: the energy space $L^2([0,T] \times \R, dt \otimes \nu)$ and the difference operator $D_J$ are the correct objects, and the modified Leibniz rule is a theorem, not an anomaly.
\end{remark}

The PDE perspective extends naturally to complex Gaussian processes, connecting to the holomorphic theory of \cite{PaperV}.

\subsection{Complex Brownian motion}

For complex Brownian motion $Z_t = X_t + iY_t$ with $X, Y$ independent standard Brownian motions, the generator acts on $\mathbb{R}^2 \cong \mathbb{C}$:
\begin{equation}
\mathcal{L} = \frac{1}{2}(\partial_x^2 + \partial_y^2) = \frac{1}{2}\Delta_{\mathbb{R}^2} = 2\partial_z\partial_{\bar{z}}.
\end{equation}

\begin{proposition}[Holomorphic Feynman--Kac]
For $f : \mathbb{C} \to \mathbb{C}$ analytic, the function $u(t,z) = \E[f(Z_T) \mid Z_t = z]$ satisfies:
\begin{enumerate}
\item[(a)] The heat equation $\partial_t u + \frac{1}{2}\Delta u = 0$.
\item[(b)] The Cauchy--Riemann equation $\partial_{\bar{z}} u = 0$ (analyticity in $z$).
\item[(c)] Therefore $\partial_t u = 0$: the conditional expectation is constant in $t$.
\end{enumerate}
\end{proposition}

\begin{proof}
Part (a) follows from Theorem~\ref{thm:main-bm} applied to the $\mathbb{R}^2$-valued process $(X_t, Y_t)$.

For (b): since $f$ is analytic and $Z_t$ is a martingale, $f(Z_t)$ is a local martingale (by It\^o's formula with $\partial_{\bar{z}}f = 0$). The conditional expectation $u(t,z) = \E[f(Z_T) \mid Z_t = z]$ inherits analyticity in $z$.

Part (c) follows from $\Delta u = 4\partial_z\partial_{\bar{z}} u = 0$ when $\partial_{\bar{z}} u = 0$.
\end{proof}

This recovers the classical result that analytic functions of Brownian motion are martingales, now understood through the Feynman--Kac lens.

\subsection{Complex fBM}

For complex fBM $Z^H_t = B^{H,(1)}_t + iB^{H,(2)}_t$ with independent components:
\begin{equation}
\partial_t u = Ht^{2H-1}\Delta u = 4Ht^{2H-1}\partial_z\partial_{\bar{z}} u.
\end{equation}

The holomorphic/anti-holomorphic structure persists: for analytic $f$, we still have $\partial_{\bar{z}} u = 0$, giving $\partial_t u = 0$. However, the energy identity of \cite[Proposition 7.6]{PaperV} shows that the holomorphic and anti-holomorphic components are \emph{not orthogonal} for $H \neq 1/2$, reflecting the non-Markovian correlation structure.

%==============================================================================
\section{The General Principle}\label{sec:general}
%==============================================================================

We now formalize the pattern observed in the preceding sections.

\begin{theorem}[Energy space determines PDE structure]\label{thm:general-principle}
Let $X$ be a stochastic process with operator triple $(D_X, \delta_X, \Pi)$ on energy space $\Hi$. Then:
\begin{enumerate}
\item[(a)] The intrinsic bracket of $X$ satisfies $d\langle X \rangle^{(D,\Pi)}_t = \|\Pi D_X X_t\|^2_\Hi \, dt$.
\item[(b)] For Gaussian $X$ with covariance $R$, this equals $\frac{d}{dt}R(t,t)\,dt$.
\item[(c)] For jump processes with L\'evy measure $\nu$, the generator includes the nonlocal term $\int(f(x+z) - f(x) - zf'(x)\mathbf{1}_{|z|\leq 1})\,\nu(dz)$.
\item[(d)] For compositions $Z = X \circ Y$, the energy space is a tensor product and the generator may be higher-order.
\end{enumerate}
\end{theorem}

\begin{proof}
Parts (a)--(b): For Gaussian processes, $\|k_t\|^2_\Hi = R(t,t)$ by the reproducing property.

Part (c): For jump processes, the difference operator $D_J$ and energy space $L^2(dt \otimes \nu)$ produce the L\'evy--Khintchine generator (Section~\ref{sec:levy}).

Part (d): For compositions, the chain rule applies at each level (Section~\ref{sec:higher}).
\end{proof}

\begin{corollary}[Classification by energy space structure]
The PDE type is determined by the energy space. When the energy space norm involves first derivatives ($\|h'\|_{L^2}$), we get second-order differential operators. When it involves second derivatives ($\|h''\|_{L^2}$), we get time-dependent coefficients. Tensor product energy spaces from compositions yield higher-order operators. Energy spaces of the form $L^2(dt \otimes \nu)$ from jump processes yield nonlocal integral operators, with power-law L\'evy measures producing fractional Laplacians.
\end{corollary}

%==============================================================================
\section{Conclusion}
%==============================================================================

We have established a comprehensive connection between energy space structure and PDE generators across Gaussian, non-Gaussian, and jump processes.

For Brownian motion, the Cameron--Martin norm $\|h'\|_{L^2}$ produces the classical Laplacian $\frac{1}{2}\Delta$. Fractional Brownian motion has a more complex RKHS $\Hi_H$ that encodes memory, producing a time-dependent coefficient $Ht^{2H-1}$ rather than a fractional spatial operator. Integrated Brownian motion, with RKHS norm $\|h''\|_{L^2}$, gives the time-dependent generator $\frac{t^2}{2}\partial_x^2$.

For iterated Brownian motion $Z_t = B^{(1)}(|B^{(2)}_t|)$, the tensor product energy space $\Hi^{(1)} \otimes \Hi^{(2)}$ and the compositional structure of the operator derivative produce the biharmonic generator $\frac{1}{8}\partial_x^4$---a new derivation of the Allouba--Zheng result.

For jump processes, the difference operator and energy space $L^2(dt \otimes \nu)$ produce nonlocal generators. The Poisson process gives $\lambda(f(n+1) - f(n))$; $\alpha$-stable processes give the fractional Laplacian $(-\Delta)^{\alpha/2}$.

The structural principle is: continuous Gaussian processes produce differential operators with coefficients determined by $\frac{1}{2}\frac{d}{dt}R(t,t)$; non-Markov Gaussian processes encode memory through time-dependent coefficients; compositions yield higher-order operators; and jumps yield nonlocal operators with kernels determined by the L\'evy measure.

The key clarification: \textbf{fractional Laplacians arise from jumps, not memory}. FBM has long-range dependence but produces a time-dependent heat equation, not a fractional PDE. Stable processes have independent increments but produce the fractional Laplacian through their jump structure.

\textbf{Open directions:} Path-dependent PDEs for non-Markovian processes; rough PDE structure for $H < 1/4$; stochastic PDEs where both the noise and the spatial operator arise from energy space structure.

\section*{Acknowledgments}

AI tools were used as interactive assistants for drafting and mathematical exploration. The author is responsible for all content.

\begin{thebibliography}{99}

\bibitem{PaperI}
R.~Fontes.
\newblock Stochastic Calculus as Operator Factorization: An Operator-Covariant Derivative and Unified Clark-Ocone--It\^o Representations.
\newblock arXiv:2601.09976, 2026.

\bibitem{PaperII}
R.~Fontes.
\newblock Stochastic Calculus for Rough Fractional Brownian Motion via Operator Factorization.
\newblock arXiv:2601.09967, 2026.

\bibitem{PaperIII}
R.~Fontes.
\newblock An It\^o Formula via Predictable Projection for Non-Semimartingale Processes.
\newblock arXiv:2601.10359, 2026.

\bibitem{PaperIV}
R.~Fontes.
\newblock The Operator Derivative: Intrinsic Properties, Structural Conditions, and Process-Dependent Extensions.
\newblock Preprint, 2026.

\bibitem{PaperV}
R.~Fontes.
\newblock The Stochastic Cauchy Operator: Operator Factorization Meets Complex Analysis.
\newblock Preprint, 2026.

\bibitem{Burdzy1993}
K.~Burdzy.
\newblock Some path properties of iterated {B}rownian motion.
\newblock In \emph{Seminar on Stochastic Processes 1992}, Progress in Probability, vol.~33, pp.~67--87. Birkh\"auser, 1993.

\bibitem{AlloubaZheng2001}
H.~Allouba and W.~Zheng.
\newblock Brownian-time processes: The {PDE} connection and the half-derivative generator.
\newblock \emph{Ann.\ Probab.}, 29(4):1780--1795, 2001.

\bibitem{Nualart2006}
D.~Nualart.
\newblock \emph{The Malliavin Calculus and Related Topics}, 2nd ed.
\newblock Springer, 2006.

\bibitem{Decreusefond1999}
L.~Decreusefond and A.S.~\"Ust\"unel.
\newblock Stochastic analysis of the fractional {B}rownian motion.
\newblock \emph{Potential Analysis}, 10:177--214, 1999.

\bibitem{DiNunno2009}
G.~Di Nunno, B.~{\O}ksendal, and F.~Proske.
\newblock \emph{Malliavin Calculus for L\'evy Processes with Applications to Finance}.
\newblock Springer, 2009.

\bibitem{Sato1999}
K.~Sato.
\newblock \emph{L\'evy Processes and Infinitely Divisible Distributions}.
\newblock Cambridge University Press, 1999.

\end{thebibliography}

\end{document}
